\documentclass[14pt,aspectratio=1610]{beamer}
%\documentclass[14pt,aspectratio=1610,handout]{beamer}
\usepackage{csu-theme}
\usepackage{forest}
\usepackage{booktabs}

\title{\Huge{} Quick Sort}
\subtitle{CSCI 315: Data Structure Analysis\\Chapter 18}
%\author{Sean T. Hayes, Ph.D.}
\institute{Charleston Southern University}

%% Comment out date to use default (today)
% \date{February 18, 2025}
\subject{Software Programming}
%\keywords{}

\setlength{\parskip}{0.5em}

\usepackage{arrays}

\tikzset{
  pointer/.append style = {
    minimum height=4mm,
    minimum width=4mm,
    node distance=1pt
  },
  label of pointer/.append style = {
    font=\ttfamily\small
  },
  node of array/.append style = {
     minimum height=8mm,
     minimum width=10.8mm,
     font=\large
  },
  label of array el/.append style={
    % rotate=-67.5,
    % anchor=west,
    % align=right,
    % yshift=-3mm,
    font=\sffamily,
    % rotate=-12,
  },
}
\begin{document}

\begin{frame}
  \titlepage%
\end{frame}

\section{Best Performance}

\begin{frame}{Lower Bound on Comparison-Based Sorting}
\center\Large
Can we do better than \(O(n^2)\) for sorting algorithms?\vspace{1em}

What is the best we can do?
\end{frame}

\begin{frame}{Comparison Tree}
\large
\begin{itemize}
\item
  \emph{Comparison tree}: graph used to trace the execution of a comparison-based algorithm.
  \begin{itemize}
  \item
    Let $L$ be a list of $n$ distinct elements, where \(n > 0\) \\
    For any $j$ and $k$, where \( 0 \le j < n \),  \( 0 \le k < n \),\\
    either \(L_{j} \le L_{k}\)  or  \(L_{j} > L_{k}\).
  \end{itemize}
\item
  \emph{Binary tree}: each comparison has two outcomes.
\end{itemize}
\end{frame}

\begin{frame}{Terminology of Comparison Trees}
\begin{columns}
  \column{0.66\textwidth}
\begin{itemize}

\item
  \emph{Node}: represents a comparison
  \begin{itemize}
  \item
    Labeled as $j:k$ (comparison of $L_{j}$ with $L_{k}$)
  \item
    If $L_{j} < L_{k}$, follow the left branch; otherwise, follow the right branch.
  \end{itemize}
\item
  \emph{Leaf}: represents the final ordering of the nodes
\item
  \emph{Root}: the top node
\item
  \emph{Branch}: line that connects two nodes
\item
  \emph{Path}: sequence of branches from one node to another
\end{itemize}
\column{0.33\textwidth}
\begin{center}
  \begin{forest}
  for tree = {
    circle, draw,
    inner sep=3pt,
    color=CSUblue,
    fill=white,
    fit=tight,
    edge=->,
    edge=very thick
  }
    [A
      [B
        [D
          [G]
          [, phantom]
        ]
        [E]
      ]
      [C
        [F]
        [, phantom]
      ]
    ]
  \end{forest}
\end{center}
\end{columns}
\end{frame}

\begin{frame}{Comparison Tree for Sorting 3 Items}
\center{}
\begin{forest}
for tree = {
  circle, draw,
  color=CSUblue,
  fill=white,
  %fit=band,
  l sep=2.6cm,
  s sep+=3mm,
  inner sep=3pt,
  edge=->,
  edge=very thick,
  grow=east,
  for leaves={rectangle, fill=CSUgold, inner sep=4pt},
  font=\small
}
[{$L_0:L_1$} ,name=root,
    [{$L_1:L_2$}, edge label = {node [midway,below,sloped] {$L_0 \le L_1$} }
      [{$\mathbf{1, 2, 3}$}, edge label = {node [midway,below,sloped] {$L_1 \le L_2$} }]
      [{$L_0:L_2$}, edge label = {node [midway,above,sloped] {$L_1 > L_2$} }
        [{$\mathbf{1, 3, 2}$}, edge label = {node [midway,below,sloped] {$L_0 \le L_2$} }]
        [{$\mathbf{3, 1, 2}$}, edge label = {node [midway,above,sloped] {$L_0 > L_2$} }]
      ]
    ]
    [{$L_1:L_2$}, edge label = {node [midway,above,sloped] {$L_0 > L_1$} }
      [{$L_0:L_2$}, edge label = {node [midway,below,sloped] {$L_1 \le L_2$} }
        [{$\mathbf{2, 1, 3}$}, edge label = {node [midway,below,sloped] {$L_0 \le L_2$} }]
        [{$\mathbf{2, 3, 1}$}, edge label = {node [midway,above,sloped] {$L_0 > L_2$} }]
      ]
      [{$\mathbf{3, 2, 1}$}, edge label = {node [midway,above,sloped] {$L_1 > L_2$} }]
    ]
]
\end{forest}
\end{frame}

\begin{frame}{Comparison Tree}
\begin{itemize}
\item<+->
  A unique permutation of the elements of $L$ is associated with each root-to-leaf path, because the sort algorithm only compares and moves the data.
\item<+->
  For a list of $n$ elements, $n > 0$, there are $n!$ different permutations.
  \begin{itemize}
  \item
    Any of these might be the correct ordering of $L$.
  \end{itemize}
\item<+->
  Thus, the tree must have \emph{at least} $n!$ leaves.
\end{itemize}
\end{frame}

\begin{frame}{Lower Bound on Comparison-Based Sorting}
\large
Theorem: Any algorithm that sorts a $n$ elements\\ (by comparison of only the keys) makes \emph{at least} $O(n\log_2n)$ key comparisons in its worst case.\vspace{1em}

\pause{}
We won't prove this theorem, but unlike the previous algorithms discussed, there are algorithms that, on average, are $O(n\log_2n)$.
\end{frame}

\section{Quick Sort}

\begin{frame}{Introduction to Quick Sort}
\emph{Quick sort}: uses the divide-and-conquer technique.

\begin{itemize}
\item
  The list is \textit{partitioned} into two sub-lists.
\item
  Each sub-list is then sorted.
\item
  Sorted sub-lists are combined into one list in such a way that the combined
  list is sorted.
\item
  All the sorting work occurs when partitioning the list.
\end{itemize}
\end{frame}

\begin{frame}{Pivot Element}
Choose an element to split the list \emph{pivot element}.
\begin{itemize}
\item
  Move elements less than the pivot value to the lower sub-list.
\item
  Move elements greater than the pivot value to the upper sub-list.
\item
  The pivot can be chosen in several ways.\\
  Ideally, the pivot will divide the list into two sub-lists of near-equal size.
\end{itemize}
\end{frame}

\section{Quick Sort: Example}

\begin{frame}[t]{Quick Sort Example}
  \vspace{2.1em}
  \centerline{%
  \begin{tikzpicture}[]
    \ArrayLabel{}
    \alt<1>{
      \ArrayNode{52}
    }{
      \ArrayNode[MyGreen]{52}
      \ArrayElLabel[val]{pivot}
    }
    \ArrayList{55, 87, 13, 78, 96, 32, 48, 22, 11, 58, 66, 88, 45}

    \ArrayGroup{el52}{el45}{Unsorted List}
  \end{tikzpicture}%
  }

  \vspace{0.7em}
  \centerline{%
  \begin{tikzpicture}
    \ArrayLabel{}
    \ArrayList[CSUgold]{45, 13, 32, 48, 22, 11}
  
    \ArrayNode[MyGreen]{52}
    \ArrayElLabel[val]{pivot}
  
    \ArrayList[MyPurple]{55, 78, 96, 58, 66, 88, 87}
  
    \ArrayGroup{el45}{el11}{Lower}
    \ArrayGroup{el55}{el87}{Upper}
  \end{tikzpicture}%
  }
\end{frame}

\begin{frame}[b]{Quick Sort Example}
\centerline{%
\begin{tikzpicture}[]
  \ArrayLabel{}
  \ArrayNode[MyGreen]{52}
  \ArrayElLabel<-12mm><0mm>[val]{pivot}
  \onslide<-4>{
      \ArrayElLabel<-8mm><8mm>[val]{small}
  }
  \alt<1>{
    \ArrayNode{55}
  }{
    \alt<-5>{
      \alt<-4>{
        \ArrayNode[MyPurple]{55}
      }{
        \ArrayNode[CSUgold]{55}
      }

    }{
      \ArrayNode[CSUgold]{13}
    }
    \onslide<5->{
      \ArrayElLabel<-8mm><0mm>[val]{small}
    }
  }

  \alt<-2>{
    \ArrayNode{87}
  }{
    \ArrayNode[MyPurple]{87}
  }

  \alt<-5>{
    \ArrayNode{13}
  }{
    \ArrayNode{55}
  }
  \ArrayList{78, 96, 32, 48, 22, 11, 58, 66, 88, 45}

  \onslide<1>{
    \ArrayGroup{el52}{el45}{Unsorted List}
  }
  \alt<1>{
    \ArrayElLabel<-4mm><4mm>[el55]{index}
  }{
    \alt<-2>{
      \ArrayElLabel[el87]{index}
      \ArrayGroup{el55}{el55}{Upper}
    }{

      \alt<-4>{
        \ArrayElLabel<-8mm><0mm>[el13]{index}
        \ArrayGroup{el55}{el87}{Upper}
      }{
        \alt<-5>{
          \ArrayElLabel<-8mm><0mm>[el13]{index}
          \ArrayGroup{el55}{el55}{Lower}
        }{
          \ArrayElLabel<-8mm><0mm>[el55]{index}
          \ArrayGroup{el13}{el13}{Lower}
        }
        \ArrayGroup<16pt>{el87}{el87}{Upper}
      }
    }
  }
  \only<4-5>{
    \ArraySwap{el55}{el13}
  }
  \only<6>{
    \ArraySwap{el13}{el55}
  }
  \onslide<7>{}

  % reserve the space for this in a later slide.
  \invisible{%
    \ArraySwap{el32}{el45}%
  }
\end{tikzpicture}%
}
\vspace{3.5em}
\end{frame}

\begin{frame}[b]{Quick Sort Example}

\centerline{%
\begin{tikzpicture}
  \ArrayLabel{}
  \ArrayNode[MyGreen]{52}
  \ArrayElLabel<-12mm><0mm>[val]{pivot}
  \ArrayNode[CSUgold]{13}
  \ArrayElLabel<-8mm><0mm>[val]{small}
  \ArrayList[MyPurple]{87, 55}
  \alt<1>{
    \ArrayNode[]{78}
  } {
    \ArrayNode[MyPurple]{78}
  }

  \alt<-2>{
    \ArrayNode{96}
  }{
    \ArrayNode[MyPurple]{96}
  }
  \ArrayList{32, 48, 22, 11, 58, 66, 88, 45}

  \ArrayGroup{el13}{el13}{Lower}
  \alt<1> {
    \ArrayElLabel[el78]{index}

    \ArrayGroup{el87}{el55}{Upper}
  }{
    \alt<-2>{
      \ArrayElLabel[el96]{index}
      \ArrayGroup{el87}{el78}{Upper}
    }{
      \ArrayElLabel[el32]{index}
      \ArrayGroup{el87}{el96}{Upper}
    }
  }
  % reserve the space for this in a later slide.
  \invisible{%
    \ArraySwap{el32}{el45}%
  }
  \onslide<3>{}
\end{tikzpicture}%
}
\vspace{3.5em}
\end{frame}

\begin{frame}[b]{Quick Sort Example}
\centerline{%
\begin{tikzpicture}
  \ArrayLabel{}
  \ArrayNode[MyGreen]{52}
  \ArrayElLabel<-12mm><0mm>[val]{pivot}
  \ArrayNode[CSUgold]{13}

  \alt<1>{
    \ArrayElLabel<-8mm><0mm>[val]{small}
    \ArrayNode[MyPurple]{87}
  }{
    \ArrayNode[CSUgold]{87}
    \ArrayElLabel{small}
  }

  \ArrayList[MyPurple]{55, 78, 96}
  \ArrayNode[]{32}
  \ArrayList{48, 22, 11, 58, 66, 88, 45}

  \ArrayElLabel[el32]{index}


  \ArraySwap{el87}{el32}

  \alt<1> {
    \ArrayGroup{el13}{el13}{Lower}
    \ArrayGroup{el87}{el96}{Upper}
  }{
    \ArrayGroup{el13}{el87}{Lower}
    \ArrayGroup{el55}{el96}{Upper}
  }

  % reserve the space for this in a later slide.
  \invisible{%
    \ArraySwap{el32}{el45}%
  }

  \onslide<2>{}
\end{tikzpicture}%
}
\vspace{3.5em}
\end{frame}


\begin{frame}[b]{Quick Sort Example}
\centerline{%
\begin{tikzpicture}
  \ArrayLabel{}
  \ArrayNode[MyGreen]{52}
  \ArrayElLabel<-12mm><0mm>[val]{pivot}
  \ArrayList[CSUgold]{13, 32}
  \ArrayElLabel[val]{small}

  \ArrayList[MyPurple]{55, 78, 96}
  \alt<-2>{
    \ArrayNode{87}
  }{
    \ArrayNode[MyPurple]{87}
  }
  \ArrayList{48, 22, 11, 58, 66, 88, 45}

  \ArrayGroup{el13}{el32}{Lower}
  \alt<-2>{
    \ArrayElLabel[el87]{index}
    \ArrayGroup{el55}{el96}{Upper}
  }{
    \ArrayElLabel[el48]{index}
    \ArrayGroup{el55}{el87}{Upper}
  }

  \onslide<1> {
    \ArraySwap{el32}{el87}
  }

  % reserve the space for this in a later slide.
  \invisible{%
    \ArraySwap{el87}{el45}%
  }

  \onslide<3> {}
\end{tikzpicture}%
}
\vspace{3.5em}
\end{frame}

\begin{frame}[b]{Quick Sort Example}
\centerline{%
\begin{tikzpicture}
  \ArrayLabel{}
  \ArrayNode[MyGreen]{52}
  \ArrayElLabel<-12mm><0mm>[val]{pivot}
  \ArrayList[CSUgold]{13, 32}

  \alt<1>{
    \ArrayNode[MyPurple]{55}
  }{
    \alt<-2>{
      \ArrayNode[CSUgold]{55}
    }{
      \ArrayNode[CSUgold]{48}
    }
  }
  \ArrayList[MyPurple]{78, 96, 87}
  \alt<-2>{
    \ArrayNode{48}
  }{
    \alt<-4>{
      \ArrayNode{55}
    }{
      \ArrayNode[MyPurple]{55}
    }
  }
  \ArrayList{22, 11, 58, 66, 88, 45}

  \onslide<1-2> {
    \ArraySwap{el55}{el48}
  }
  \onslide<3> {
    \ArraySwap{el48}{el55}
  }

  \alt<-2>{
    \ArrayElLabel[el48]{index}
  }{
    \alt<-4>{
      \ArrayElLabel[el55]{index}
    }{
      \ArrayElLabel[el22]{index}
    }
  }

  \alt<1> {
    \ArrayElLabel[el32]{small}
    \ArrayGroup{el13}{el32}{Lower}
    \ArrayGroup{el55}{el87}{Upper}
  }{
    \alt<-2>{
      \ArrayElLabel[el55]{small}
      \ArrayGroup{el13}{el55}{Lower}
    }{
      \ArrayElLabel[el48]{small}
      \ArrayGroup{el13}{el48}{Lower}
    }
    \alt<-4>{
      \ArrayGroup{el78}{el87}{Upper}
    }{
      \ArrayGroup{el78}{el55}{Upper}
    }
  }

  % reserve the space for this in a later slide.
  \invisible{%
    \ArraySwap{el87}{el45}%
  }

  \onslide<5>{}
\end{tikzpicture}%
}
\vspace{3.5em}
\end{frame}

\begin{frame}[b]{Quick Sort Example}
\centerline{%
\begin{tikzpicture}
  \ArrayLabel{}
  \ArrayNode[MyGreen]{52}
  \ArrayElLabel<-12mm><0mm>[val]{pivot}
  \ArrayList[CSUgold]{13, 32, 48}

  \alt<1>{
    \ArrayElLabel[val]{small}
    \ArrayNode[MyPurple]{78}
  }{
    \alt<-2>{
      \ArrayNode[CSUgold]{78}
    }{
      \ArrayNode[CSUgold]{22}
    }
    \ArrayElLabel[val]{small}
  }

  \ArrayList[MyPurple]{96, 87, 55}

  \alt<-2>{
    \ArrayNode{22}
  }{
    \alt<-4>{
      \ArrayNode{78}
    } {
      \ArrayNode[MyPurple]{78}
    }
  }

  \alt<-4>{
    \ArrayElLabel{index}
    \ArrayNode{11}
  }{
    \ArrayNode{11}
    \ArrayElLabel{index}
  }

  \ArrayList{58, 66, 88, 45}

  \onslide<1-2> {
    \ArraySwap{el78}{el22}
  }
  \onslide<3> {
    \ArraySwap{el22}{el78}
  }

  \alt<1> {
    \ArrayGroup{el13}{el48}{Lower}
    \ArrayGroup{el78}{el55}{Upper}
  }{
    \alt<-2>{
      \ArrayGroup{el13}{el78}{Lower}
    }{
      \ArrayGroup{el13}{el22}{Lower}
    }
    \alt<-4>{
      \ArrayGroup{el96}{el55}{Upper}
    }{
      \ArrayGroup{el96}{el78}{Upper}
    }
  }

  % reserve the space for this in a later slide.
  \invisible{%
    \ArraySwap{el87}{el45}%
  }

  \onslide<5>{}
\end{tikzpicture}%
}
\vspace{3.5em}
\end{frame}

\begin{frame}[b]{Quick Sort Example}
\centerline{%
\begin{tikzpicture}
  \ArrayLabel{}
  \ArrayNode[MyGreen]{52}
  \ArrayElLabel<-12mm><0mm>[val]{pivot}
  \ArrayList[CSUgold]{13, 32, 48, 22}

  \alt<1>{
    \ArrayElLabel[val]{small}
    \ArrayNode[MyPurple]{96}
  }{
    \alt<-2>{
      \ArrayNode[CSUgold]{96}
    }{
      \ArrayNode[CSUgold]{11}
    }
    \ArrayElLabel[val]{small}
  }

  \ArrayList[MyPurple]{87, 55, 78}

  \alt<-2>{
    \ArrayNode{11}
  }{
    \alt<-4>{
      \ArrayNode{96}
    } {
      \ArrayNode[MyPurple]{96}
    }
  }

  \alt<-4>{
    \ArrayElLabel{index}
    \ArrayNode{58}
  }{
    \ArrayNode{58}
    \ArrayElLabel{index}
  }

  \ArrayList{66, 88, 45}

  \onslide<1-2> {
    \ArraySwap{el96}{el11}
  }
  \onslide<3> {
    \ArraySwap{el11}{el96}
  }

  \alt<1> {
    \ArrayGroup{el13}{el22}{Lower}
    \ArrayGroup{el96}{el78}{Upper}
  }{
    \alt<-2>{
      \ArrayGroup{el13}{el96}{Lower}
    }{
      \ArrayGroup{el13}{el11}{Lower}
    }
    \alt<-4>{
      \ArrayGroup{el87}{el78}{Upper}
    }{
      \ArrayGroup{el87}{el96}{Upper}
    }
  }

  % reserve the space for this in a later slide.
  \invisible{%
    \ArraySwap{el87}{el45}%
  }

  \onslide<5>{}
\end{tikzpicture}%
}
\vspace{3.5em}
\end{frame}


\begin{frame}[b]{Quick Sort Example}
\centerline{%
\begin{tikzpicture}
  \ArrayLabel{}
  \ArrayNode[MyGreen]{52}
  \ArrayElLabel<-12mm><0mm>[val]{pivot}
  \ArrayList[CSUgold]{13, 32, 48, 22, 11}
  \ArrayElLabel[val]{small}

  \ArrayList[MyPurple]{87, 55, 78, 96, 58}
  \alt<1> {
    \ArrayNode[]{66}
  } {
    % \ArrayNode[MyPurple]{66}
  }
  \ArrayList{88, 45}

  \ArrayElLabel[el66]{index}

  \onslide<1> {
    \ArrayGroup{el13}{el11}{Lower}
    \ArrayGroup{el87}{el58}{Upper}
  }
  % \onslide<2> {
  %   \ArrayGroup{el13}{el11}{Lower}
  %   \ArrayGroup{el87}{el66}{Upper}
  % }

  % reserve the space for this in a later slide.
  \invisible{%
    \ArraySwap{el87}{el45}%
  }

\end{tikzpicture}%
}
\vspace{3.5em}
\end{frame}

\begin{frame}[b]{Quick Sort Example}
\centerline{%
\begin{tikzpicture}
  \ArrayLabel{}
  \ArrayNode[MyGreen]{52}
  \ArrayElLabel<-12mm><0mm>[val]{pivot}
  \ArrayList[CSUgold]{13, 32, 48, 22, 11}
  \ArrayElLabel{small}

  \ArrayList[MyPurple]{87, 55, 78, 96, 58, 66}
  \alt<1> {
    \ArrayNode[]{88}
    \ArrayElLabel[el88]{index}
  } {
    \ArrayNode[MyPurple]{88}
  }
  \ArrayList{45}

  \onslide<2->{
    \ArrayElLabel[el45]{index}
  }

  \ArrayGroup{el13}{el11}{Lower}
  \alt<1> {
    \ArrayGroup{el87}{el66}{Upper}
  }{
    \ArrayGroup{el87}{el88}{Upper}
  }

  % reserve the space for this in a later slide.
  \invisible{%
    \ArraySwap{el87}{el45}%
  }
\end{tikzpicture}%
}
\vspace{3.5em}
\end{frame}

\begin{frame}[b]{Quick Sort Example}
\centerline{%
\begin{tikzpicture}
  \ArrayLabel{}
  \ArrayNode[MyGreen]{52}
  \ArrayElLabel<-12mm><0mm>[val]{pivot}
  \ArrayList[CSUgold]{13, 32, 48, 22, 11}

  \alt<1>{
    \ArrayElLabel{small}
    \ArrayNode[MyPurple]{87}
  }{
    \alt<-2>{
      \ArrayNode[CSUgold]{87}
    }{
      \ArrayNode[CSUgold]{45}
    }
    \ArrayElLabel[val]{small}
  }

  \ArrayList[MyPurple]{55, 78, 96, 58, 66, 88}



  \alt<-2>{
    \ArrayNode[]{45}
  }{
    \alt<-4>{
      \ArrayNode[]{87}
    } {
      \ArrayNode[MyPurple]{87}
    }
  }

  \ArrayElLabel{index}


  \onslide<1-2> {
    \ArraySwap{el87}{el45}
  }
  \onslide<3> {
    \ArraySwap{el45}{el87}
  }

  %\onslide<1> {

  %}

  \alt<1> {
    \ArrayGroup{el13}{el11}{Lower}
    \ArrayGroup{el87}{el88}{Upper}
  }{
    \alt<-2>{
      \ArrayGroup{el13}{el87}{Lower}
    }{
      \ArrayGroup{el13}{el45}{Lower}
    }
    \alt<-4>{
      \ArrayGroup{el55}{el88}{Upper}
    }{
      \ArrayGroup{el55}{el87}{Upper}
    }
  }
  % reserve the space for this in a later slide.
  \invisible{%
    \ArraySwap{el87}{el45}%
  }
\end{tikzpicture}%
}
\vspace{3.5em}
\end{frame}

\begin{frame}[b]{Quick Sort Example}
\centerline{%
\begin{tikzpicture}
  \ArrayLabel{}
  \ArrayNode[MyGreen]{52}
  \ArrayElLabel<-12mm><0mm>[val]{pivot}
  \ArrayList[CSUgold]{13, 32, 48, 22, 11, 45}
  \ArrayElLabel[val]{small}

  \ArrayList[MyPurple]{55, 78, 96, 58, 66, 88, 87}

  \onslide<2> {
    \ArraySwap{el52}{el45}
  }

  \ArrayGroup{el13}{el45}{Lower}
  \ArrayGroup{el55}{el87}{Upper}

  % reserve the space for this in a previous slide.
  \invisible{%
    \ArraySwap{el45}{el87}%
  }
\end{tikzpicture}%
}
\vspace{3.5em}
\end{frame}


\begin{frame}[b]{Quick Sort Example}
\centerline{%
\begin{tikzpicture}
  \ArrayLabel{}
  \ArrayList[CSUgold]{45, 13, 32, 48, 22, 11}

  \ArrayNode[MyGreen]{52}
  \ArrayElLabel[val]{pivot}

  \ArrayList[MyPurple]{55, 78, 96, 58, 66, 88, 87}

  \onslide<1> {
    \ArraySwap{el45}{el52}
  }

  \ArrayGroup{el45}{el11}{Lower}
  \ArrayGroup{el55}{el87}{Upper}

  % reserve the space for this in a previous slide.
  \invisible{%
    \ArraySwap{el52}{el87}%
  }
\end{tikzpicture}%
}
\pause\vspace{3.5em}
\end{frame}

% \section{Quick Sort: Implementation}

% \begin{frame}[fragile]{Quick Sort Recursive Implementation}
% \begin{lstlisting}[basicstyle=\small\ttfamily]
% // Array-based implementation
% template <typename Type>
% void quickSort(Type arr[], int size) {

%   if (size > 1) {
%     const int PIVOT_INDEX {partition(arr, size)};
%     const int UPPER_FIRST {INDEX + 1};

%     quickSort(arr, PIVOT_INDEX); // lower half
%     quickSort(arr + UPPER_FIRST, size - UPPER_FIRST);
%   }
% }
% \end{lstlisting}
% \end{frame}

% \begin{frame}[fragile]{Quick Sort Recursive Implementation}
% \begin{lstlisting}[basicstyle=\small\ttfamily]
% template <typename Type>
% int partition(Type arr[], int size) {
%   int smallIndex = 0;

%   for (int index {1}; index < size; ++index) {
%     if (arr[index] < arr[0]) { // less than pivot
%       ++smallIndex;
%       swap(arr[index], arr[smallIndex]);
%     }
%   }

%   // put the pivot in its place
%   swap(arr[0], arr[smallIndex]);
%   return smallIndex;
% }
% \end{lstlisting}
% \end{frame}

\section{Quick Sort: Analysis}

\begin{frame}{Analysis of Quick Sort}
\begin{tabular}{r l l}
\toprule
\textbf{Algorithm}        & \textbf{Comparisons} & \textbf{Swaps} \\
\midrule
  Bubble Sort    & $O(n^2)$ & $O(n^2)$ \\
  Selection Sort & $O(n^2)$ & $O(n)$ \\
  Array-Based Insertion Sort & $O(n^2)$ & $O(n^2)$ \\
  Linked-List Insertion Sort & $O(n^2)$ & $O(n)$ \\
  Quick Sort: Average Case& \onslide<2->{$O(n \log_2{n})$} &  \onslide<2->{$O(n \log_2{n})$} \\
  Quick Sort: Worst Case &  \onslide<3->{$O(n^2)$} &  \onslide<3->{$O(n^2)$} \\
\bottomrule
\end{tabular}

\begin{tabular}{r c c}
  \toprule
  \textbf{Quick Sort}        & \textbf{Comparisons } & \textbf{Swaps} \\
  \midrule
    Average Case & \onslide<2->{$(1.39) n \log_2{n}$} & \onslide<2->{$(0.69) n \log_2{n}$} \\
    Worst Case   &  \onslide<3->{$\frac{n^2 - n}{2}$} & \onslide<3->{$\frac{n^2 - 3n}{2}$} \\
  \bottomrule
\end{tabular}

% \begin{itemize}
% \item<2-> Average Case:
%   \begin{itemize}
%   \item
%     Comparisons: 
%   \item
%     Swaps: 
%   \end{itemize}
% \item<3-> Worst Case:
%   \begin{itemize}
%   \item
%     Comparisons:
%   \item
%     Swaps: 
%   \end{itemize}
% \end{itemize}

\end{frame}


\begin{frame}{Worst Case of Quick Sort}
\large\center{}
  If the chosen pivot is always the smallest or largest value, Quick Sort behaves like a \emph{Selection Sort}.

\end{frame}

\begin{frame}{Quick Sort Optimizations}
\begin{itemize}
  \item Choose a better pivot.
  \begin{itemize}
    \item Median of three: first, middle, last elements.
    \item Randomly select a pivot.
  \end{itemize}
  \item Use a different algorithm for small sub-lists.
  \begin{itemize}
    \item
      Insertion sort is best for lists of 10 or fewer elements.
    \item
      The overhead of the recursive calls is greater than the cost of the insertion sort.
  \end{itemize}
  \item Use fancy partitioning techniques.
  \begin{itemize}
    \item
      Hoare's partitioning: Traverse the list from both sizes, swapping smaller elements to the left and larger elements to the right.
    \item
      Dual-Pivot partitioning: Partitions the list into three sub-lists.
  \end{itemize}
\end{itemize}
\end{frame}

\end{document}
